\section{What does not transfer}
\label{sec:limits}

\paragraph{The flow-generation condition buys reproducibility, not accuracy.} We select $C$ by asking for two things: that it stay constant along the model's trajectories,
and that it generate the flow through $F \approx B \nabla C$. The second condition is
the part that distinguishes this from a generic search for conserved quantities, and on this model it is close to
inert. We compared the recovered $C$ against 200 random degree-4 polynomials of matched complexity,
scored on the same latent. The recovered $C$ lands at the 17th to 66th percentile of that null on
pairing residual, which is mid-pack. The dissipative models' candidates land at the 100th percentile,
worse than every random polynomial, so the statistic is not vacuous. On this model, conservation does
the discriminating.

An ablation isolates what the term contributes. Holding the candidate family, data, dimension,
degree and $\alpha$ grid fixed, we select $C$ two ways: by conservation alone, taking the
lowest-ratio eigenvector, and jointly with $B$ as we do throughout. Conservation alone recovers the
energy at $0.950$ against the joint rule's $0.973$, and it edits marginally better ($-0.087$ against
$-0.077$, winning on two seeds of three). The one place the two differ is spread: the joint rule's
recovery varies by $0.008$ across seeds and conservation alone varies by $0.035$, $4.4\times$ wider.
Flow-generation makes the extraction repeatable rather than more accurate, and it is not what makes
the intervention work.

Figure~\ref{fig:ldsweep} sharpens the problem. As the extraction dimension rises from 6 to 12,
recovery improves from $0.721$ to $0.967$ while the pairing residual gets \emph{worse}, from $0.738$
to $0.876$. Recovery is already at $0.967$ by dimension 8, where the residual is better ($0.813$), so
the dimension we registered is not where recovery arrives. It is past it. We registered the
natural explanation for the anti-correlation, that the added directions carry the energy but
resist a Hamiltonian fit, and it failed: on two seeds of three those directions carry more than
twice the energy information of the first six ($2.23\times$ and $2.13\times$), but their share of the
residual (44.9\%) tracks their share of the flow (43.5\%) almost exactly (Figure~\ref{fig:lowvar}).
On the third seed the energy relation inverts ($0.91\times$), which does not affect the residual
result but does bound the claim below it. We cannot
explain the anti-correlation, and we report the pairing residual as a diagnostic rather than as a
confidence score.

\begin{figure}[t]
\centering
\begin{minipage}[t]{0.36\textwidth}
\centering
\includegraphics[width=\textwidth]{figures/fig2_ld_sweep.pdf}
\caption{Recovery and the pairing residual move in opposite directions above dimension 6, so the
residual cannot be read as a confidence score.}
\label{fig:ldsweep}
\end{minipage}\hfill
\begin{minipage}[t]{0.58\textwidth}
\centering
\includegraphics[width=\textwidth]{figures/fig3_low_variance.pdf}
\caption{Medians over three seeds. The added low-variance directions carry twice the energy
information of the top six \textbf{(a)} on two seeds of three and $0.91\times$ on the third, but on
every seed take no more than their share of the misfit \textbf{(b)}, which rules out the explanation
we had registered for the anti-correlation on the left.}
\label{fig:lowvar}
\end{minipage}
\end{figure}

\paragraph{Energy lives in the low-variance directions, on two seeds of three.} Principal
directions 7 through 12 carry $2.23\times$ and $2.13\times$ the energy information of the top six on
two seeds, and $0.91\times$ on the third, where the relation inverts. We had seen the effect on a
sharply anisotropic recurrent network, and it reproduces here on a near-isotropic stochastic model
in two models of three. We state it per seed because a median would hide a sign change. It is the
information-theoretic counterpart of the causal measurement in Section~\ref{sec:leverage}: variance
misses the directions that carry the physics, and separately misses the directions whose
displacement damages the rollout.

\paragraph{A mechanism from the GRU study does not replicate.} The GRU study showed that phase error accumulates as the time integral of the invariant error weighted by
frequency sensitivity, $\int \kappa(E_s)\,\Delta E(s)\,ds$ with $\kappa = d\ln\omega/dE$, beating a
fitted power law on held-out data by $5\times$. Here the weaker part holds: accumulated invariant
error predicts which rollouts drift most, with rank correlation rising to $0.503$ as the horizon
grows, and the integral predictor beats a fitted $t^p$ ($R^2$ $0.187$ against $0.021$). Shuffling
destroys it. The frequency weighting specifically does not separate. Under a nested test with 400
bootstrap resamples, the weighted term contributes significantly beyond the unweighted integral in 1
of 9 seed-horizon combinations, and the unweighted beyond the weighted in 0 of 9. The two predictors
are not distinguishable at $n = 512$.

We recorded this as a refutation once and were wrong. We had compared two standalone $R^2$ values
differing by $0.019$, never checked that gap against its sampling error, and let a three-seed median
choose the direction while the weighted predictor actually led in 6 of 9 combinations. The verdict is
unresolved.

\paragraph{Scope.} One physical system, one degree of freedom, one architecture, and a training run
far shorter than a standard DreamerV3 schedule. The pendulum straddles its separatrix, so 17\% of
trajectories rotate rather than oscillate and we restrict the mechanism analysis to the librating
ones. We have not tested contacts, multiple objects, actions, or any system where the frequency
varies smoothly enough to identify $\kappa$.
